\documentclass[a4paper, 12pt]{article}
\usepackage{fontspec}
\usepackage{xeCJK}
\usepackage[left=2cm, right=2cm, top=2cm, bottom=2cm]{geometry}
\usepackage{tcolorbox}
\usepackage{fancyhdr}
\usepackage{amsmath}
\pagestyle{fancy}
\lhead{高中物理}
\rhead{南昌学而思}
\cfoot{generated by Linghui Ding}
\newcommand{\defitem}[2]{
  \noindent \textsf{#1} \hspace{1em} \itshape{#2} 
}
\newcommand{\formula}[2]{
  \noindent \textbf{#2}  \[ #1 \]
}
\newtcolorbox{arcbox}{colframe=orange, arc=3mm, boxrule=1pt}
\newcommand{\noticebox}[1]{
  \begin{arcbox}
    \normalfont #1
  \end{arcbox}
}
\begin{document}
\defitem{电势}{
  计算式为
  \[\varphi = \frac W q\]
  物体在电场中某点的电势能与该点电荷量的比值，单位为伏特（V）。
}
\noticebox{
  电势本身由电荷与测验点位置决定，不因电荷与电势能而影响。
}
\defitem{电势差}{
  \[U=\varphi_a-\varphi_b\]
  两点之间的电势之差，类似于重力的高度差。\\
}
\defitem{电势能}{
  \[W=\varphi q\]
  物体在电场中某点的电势能与该点电荷量的乘积，单位为焦耳（J），类似于重力势能。\\
}
\defitem{电势能的变化}{
  \[
    \Delta W = q \Delta \varphi
  \]
  物体在电场中移动时，电势能的变化等于电荷量与电势差的乘积。
  \[
    \Delta W = W_B - W_A = q(\varphi_B - \varphi_A)
  \]
  其中 $W_A$、$W_B$ 分别为 A、B 两点的电势能，$\varphi_A$、$\varphi_B$ 分别为 A、B 两点的电势。
  \[
    W = q\varphi
  \]
  电势能等于电荷量与该点电势的乘积。
  若 $q>0$，电荷自高电势向低电势运动，电势能减少。
  若 $q<0$，电荷自低电势向高电势运动，电势能减少。
}
\noticebox{
  电场力是保守力，所以其对电荷所做的功只与初末位置有关，与路径无关。
}
\defitem{电容器}{
  存储电能的元件，由两个导体之间隔着绝缘材料构成。
  平行板电容器是由两块平行放置的导体板组成，中间夹有绝缘介质。常用于研究电容器的基本性质。\\
}
\defitem{电容}{
  描述电容器储存电能能力的物理量，定义为电容器中储存的电荷量与电压之比，单位为法拉（F）。
  平行板电容器的电容表达式为
  \[
    C = \varepsilon_0 \varepsilon_r \frac{S}{d}
  \]
  其中 $C$ 为电容，$\varepsilon_0$ 为真空介电常数，$\varepsilon_r$ 为介质的相对介电常数，$S$ 为极板面积，$d$ 为极板间距离。

  库仑定律中的系数 $k$ 与真空介电常数 $\varepsilon_0$ 的关系为
  \[
    k = \frac{1}{4\pi \varepsilon_0}
  \]
  其中 $k$ 是库仑定律中的比例常数，$\varepsilon_0$ 是真空介电常数。\\
}
\defitem{电容器的充放电}{
  电容器的充电过程是指在外加电压作用下，电容器两极板之间积累电荷的过程；放电过程是指电容
  器释放储存的电能，电荷流向外部电路的过程。只有在电容器两极板间存在电势差（即外加电源或
  电压）时，电容器才能充电。充电过程中，电流方向为正极板流入、负极板流出，直到两极板间的
  电势差等于外加电压时，充电过程结束，电流为零。\\
  \textbf{充电过程中电流和电压的变化：}
  电容器刚开始充电时，极板间电压为零，电流最大；
  随着充电进行，极板间电压逐渐升高，电流逐渐减小；
  当极板间电压等于外加电压时，电流减小到零，充电过程结束。
  电容器存在击穿电压，当电压过大时电容器中间的绝缘材料就会被击穿，使得电容器被破坏。\\
}
\defitem{电容器的能量}{
  电容器储存的电能可以用公式 $E = \frac{1}{2} C U^2=\frac12\frac {Q^2} C=\frac QU$ 
  表示，其中 $E$ 为电能，$C$ 为电容，$U$ 为电压。\\
}
\formula{E = \frac{1}{2} C U^2 = \frac{1}{2}\frac{Q^2}{C} = \frac{Q U}{2}}{电容器的能量}
\noticebox{
  电容器可以进行串并联，并联电容器，等效电容为电容器电容之和，串联电容器，等效电容规律和并联电阻相同。
}
\defitem{直线运动}{
  带电粒子在电场中受到的力是恒定的，因此其运动轨迹为直线。
  \textbf{条件：} 电场为匀强电场，带电粒子初速度与电场方向平行或反平行。
  \textbf{类型：} 匀加速直线运动（若初速度与电场方向同向或反向）。
  \textbf{加速度表达式：}
  \[
    a = \frac{F}{m} = \frac{qE}{m}
  \]
  其中 $q$ 为粒子电荷量，$E$ 为电场强度，$m$ 为粒子质量。\\
}
\defitem{类平抛运动}{
  带电粒子在电场中沿水平方向以初速度 $v_0$ 运动，同时受到垂直方向的电场力作用，形成类平
  抛运动。\\
}
\defitem{电流}{
  \textbf{电流的直接产生原因？}
  载流子\underline{定向}移动。一般在金属导体中，载流子为自由电子。
  \textbf{闭合电路的电流产生原因？}
  电源施加电压，电压在导体两端形成电场，电场力促使自由电子定向移动。
  \textbf{电流的方向？}
  正电荷移动的方向，或者负电荷移动的反方向。
  \textbf{电流的大小？}
  单位时间内通过导体单位横截面积的电荷量为电流大小。下面是电流的定义式：
  \[I=\frac Qt\]
  \textbf{恒定电流？}
  电流的大小和方向都不变的情况，并且电流仍然是一个标量，因为它没有矢量运算法则。
  \textbf{电流的决定式：}
  \[I=neSv\]
  里面的n是单位体积内自由电子数量。v是自由电子的移动速度。
}
\noticebox{
  电流决定式中，面积需要垂直于移动速度的那个投影面积。
}
\defitem{电阻}{
  导体中阻碍自由电子移动的情况。阻碍越严重，电阻就越大。
  \textbf{电阻定义式}
  \[R=\frac UI\]
  \textbf{电阻决定式}
  \[R=\rho\frac lS\]
  $l$是导体长度，$S$为导体的横截面积。
  \textbf{$\rho$电阻率影响因素？}
  金属类别：银<铜<铝<其他。
  温度，温度越大，电阻越大\\
}
\defitem{电源}{
  通过非静电力做功完成电子的定向移动循环的装置。
  电源内部，电场力做负功，电势能增加说明必须要通过其他形式给循环输入能量，这就是非静电力，
  静电力是不可能完成该循环的。\\
  \textbf{电流流向}
  在电源外部，电流从正极流向负极。在电源的内部，电流从负极流向正极。\\
  \textbf{自由电子移动方向}
  自由电子从负极移动到正极，并且在正极通过非静电力回到负极。\\
  \textbf{电动势}
  描述电源的非静电力做功能力的物理量。数值上等于电源中非静电力将1C的电荷从负极移动到正极
  所做的功。\\
  \textbf{电压与电动势}
  电压是在电源外，电场力将自由电子从负极移动到正极的能力，而电动势是电源内部，非静电力将
  电子搬回负极的能力。\\
  \textbf{电源内阻}
  电源内部的电阻，可以认为一个电源为一个内阻与理想电源进行串联。\\
}
\defitem{直线运动}{
带电粒子在电场中受到的力是恒定的，因此其运动轨迹为直线。
  \textbf{条件：} 电场为匀强电场，带电粒子初速度与电场方向平行或反平行。\\
  \textbf{类型：} 匀加速直线运动（若初速度与电场方向同向或反向）。\\
  \textbf{加速度表达式：}
    \[
      a = \frac{F}{m} = \frac{qE}{m}
    \]
    其中 $q$ 为粒子电荷量，$E$ 为电场强度，$m$ 为粒子质量。\\
}
\defitem{类平抛运动}{
  带电粒子在电场中沿水平方向以初速度 $v_0$ 运动，同时受到垂直方向的电场力作用，形成类平抛运动。
}
\defitem{电流}{
  \textbf{电流的直接产生原因？}
  载流子\underline{定向}移动。一般在金属导体中，载流子为自由电子。\\
  \textbf{闭合电路的电流产生原因？}
  电源施加电压，电压在导体两端形成电场，电场力促使自由电子定向移动。\\
  \textbf{电流的方向？}
  正电荷移动的方向，或者负电荷移动的反方向。\\
  \textbf{电流的大小？}
  单位时间内通过导体单位横截面积的电荷量为电流大小。下面是电流的定义式：
  \[I=\frac Qt\]
  \textbf{恒定电流？}
  电流的大小和方向都不变的情况，并且电流仍然是一个标量，因为它没有矢量运算法则。\\
  \textbf{电流的决定式：}
  \[I=neSv\]
  里面的n是单位体积内自由电子数量。v是自由电子的移动速度。
}
\noticebox{电流决定式中，面积需要垂直于移动速度的那个投影面积。}
\defitem{电阻}{
  导体中阻碍自由电子移动的情况。阻碍越严重，电阻就越大。\\
  \textbf{电阻定义式}
  \[R=\frac UI\]
  \textbf{电阻决定式}
  \[R=\rho\frac lS\]
  $l$是导体长度，$S$为导体的横截面积。\\
  \textbf{$\rho$电阻率影响因素？}
  金属类别：银<铜<铝<其他。
  温度，温度越大，电阻越大
}\\
\defitem{电源}{
  通过非静电力做功完成电子的定向移动循环的装置。
  电源内部，电场力做负功，电势能增加说明必须要通过其他形式给循环输入能量，这就是非静电力，
  静电力是不可能完成该循环的。\\
  \textbf{电流流向}
  在电源外部，电流从正极流向负极。在电源的内部，电流从负极流向正极。\\
  \textbf{自由电子移动方向}
  自由电子从负极移动到正极，并且在正极通过非静电力回到负极。\\
  \textbf{电动势}
  描述电源的非静电力做功能力的物理量。数值上等于电源中非静电力将1C的电荷从负极移动到正极
  所做的功。\\
  \textbf{电压与电动势}
  电压是在电源外，电场力将自由电子从负极移动到正极的能力，而电动势是电源内部，非静电力将
  电子搬回负极的能力。\\
  \textbf{电源内阻}
  电源内部的电阻，可以认为一个电源为一个内阻与理想电源进行串联。\\
}
\end{document}
